% =====================================================================
% paper/main.tex — RIVF 2026 SS3 scaffold (v2)
% Built 21/09/2026 by Ho Du Tuan Dat, per notes/2026-09-21_hung_supervisor.md
% (Priority 1, due end of 22/09). Revised after self-review against
% Hung's exact wording:
%   - "deployment and serving SUBSECTION" -> demoted from a top-level
%     \section to a \subsection inside Results (it is an experimental
%     result table, owned by Dat).
%   - "Method, threat model, and experimental setup" (one phrase, one
%     owner: Chinh) -> Threat Model folded into a \subsection of
%     Method, not a standalone \section.
%   - Limitations (not named anywhere in Hung's writing schedule) ->
%     folded into a closing \subsection of Discussion instead of its
%     own \section, to protect the hard 6-page budget he stressed
%     repeatedly ("cut to exactly six pages").
% This cuts top-level \section count from 9 to 6 + Acknowledgment.
% Structure only: every heading is final, every figure/table is a
% named placeholder tied to a specific experiment already defined in
% the plan. Fill placeholders as results land; do not invent numbers.
% =====================================================================
\documentclass[10pt,conference]{IEEEtran}
\usepackage{cite}
\usepackage{amsmath,amssymb,amsfonts}
\usepackage{graphicx}
\usepackage{textcomp}
\usepackage{xcolor}
\usepackage{booktabs}
\usepackage{hyperref}

% --- Convenience macro to flag unfilled content. Search "TODO" before
%     the 26/09 draft freeze; nothing with a TODO should remain by then.
\newcommand{\todo}[1]{\textcolor{red}{[TODO: #1]}}

\begin{document}

% ---------------------------------------------------------------------
% TITLE — placeholder. Revisit on 25/09 when Introduction is drafted
% together (Dan + Dat).
% ---------------------------------------------------------------------
\title{Backdoor Repair Under Shard-Version Skew in Distributed Edge Inference}
% NOTE: this is the project's official working title (README.md,
% Project Agreement.md). It describes the mechanism, not the attack
% framing Section 4 of the 21/09 plan calls for. Revisit only if the
% team explicitly agrees to reframe it on 24--25/09 -- do not change
% unilaterally while drafting Introduction/Abstract.

\author{
\IEEEauthorblockN{1\textsuperscript{st} Duong Ngoc Linh Dan}
\IEEEauthorblockA{Faculty of Technology\\
Van Lang University\\
Ho Chi Minh City, Vietnam\\
dan.2374802010091@vanlanguni.vn}
\and
\IEEEauthorblockN{2\textsuperscript{nd} Ho Du Tuan Dat}
\IEEEauthorblockA{Faculty of Technology\\
Van Lang University\\
Ho Chi Minh City, Vietnam\\
dat.2374802010097@vanlanguni.vn}
\and
\IEEEauthorblockN{3\textsuperscript{rd} Nguyen Minh Chinh}
\IEEEauthorblockA{Faculty of Technology\\
Van Lang University\\
Ho Chi Minh City, Vietnam\\
chinh.2275106050051@vanlanguni.vn}
% \and  % uncomment and add a 4th block here ONLY once the team
%       % confirms Hung is a listed co-author (see \todo below) --
%       % then DELETE the "Supervised by" line right below, since a
%       % listed co-author is not also captioned as "supervisor" in
%       % most venues' author blocks.
\\[4pt]
{\footnotesize Supervised by Dr. Dang Khanh Hung}
}

% \todo{Confirm with Hung and the team: is he a formal co-author
% (uncomment the 4th \IEEEauthorblockN/\IEEEauthorblockA block above,
% with his own affiliation and email, and delete the "Supervised by"
% line), or does he stay off the author list and get thanked only in
% the Acknowledgment section at the end of the paper? Do not decide
% this unilaterally -- it determines who is credited as an author on
% a submitted publication.}

\maketitle

% =====================================================================
\begin{abstract}
\todo{Write LAST, on 26/09, after every number below is frozen (23/09
freeze). Should state: (1) the setting -- distributed edge inference
with independently-signed model shards; (2) the attack -- an artifact
cache that can serve a mix of repaired and historically-valid
(pre-repair) shards; (3) the headline empirical result -- single-shard
rollback alone is weak (max ASR 17.46\% vs.\ 97.43\% for the fully
infected model) but is invisible to utility monitoring (clean accuracy
stays within a few points of normal), and recovery grows with the
number of stale shards and with how much of the repair update they
carry; (4) the defense -- epoch-manifest pinning, framed as applying
existing supply-chain controls (TUF/Uptane) to this ML-specific gap,
not as a new protocol.}
\end{abstract}

\begin{IEEEkeywords}
\todo{4--6 keywords, e.g.\ backdoor attacks, model partitioning, edge
inference, supply-chain security, version skew.}
\end{IEEEkeywords}

% =====================================================================
\section{Introduction}
\label{sec:intro}
% Owner: Dan + Dat, together, 25/09.
\todo{Introduction -- follow the attack-paper shape from
notes/2026-09-21\_hung\_supervisor.md Section 4:
\begin{itemize}
  \item Concrete deployment scenario (from Project Agreement.md,
        Section 2 "Concrete deployment example"): edge inference
        service detects a backdoor, fully repairs and re-signs the
        model, but an interrupted rollout leaves some workers on
        stale shards; per-shard signature verification does not check
        \emph{epoch} consistency across shards.
  \item State the threat model summary in one paragraph (full detail
        is in Section~\ref{sec:threat}).
  \item State contributions as a bulleted list ending the section:
    \begin{enumerate}
      \item Quantitative measurement of backdoor-repair fragility
            under shard rollback, across single-shard and multi-shard
            (\(k\)-shard) settings.
      \item A comparison of shard-selection strategies (repair-aware
            vs.\ architectural vs.\ random), showing an attacker who
            targets high-repair-share shards needs fewer stale shards
            to succeed.
      \item A layer/epoch fragility measurement contrasting a global
            20-epoch fine-tune repair against a local ANP repair.
      \item A serving-level comparison of per-shard signature checking
            vs.\ epoch-manifest pinning as a practical defense.
    \end{enumerate}
  \item Per Section 6 of the 21/09 plan: do NOT claim the team
        predicted the repair-aware result in advance -- it is reported
        as a comparison, not a confirmed hypothesis.
\end{itemize}}

% =====================================================================
\section{Method and Experimental Setup}
\label{sec:method}
% Owner: Chinh, 24/09 ("writes Method, threat model, and experimental
% setup" -- one section, one owner, per notes/2026-09-21_hung_supervisor.md
% Section 5).

\subsection{Threat Model}
\label{sec:threat}
\todo{Reuse verbatim from ``Project Agreement.md'', Section 3:
attacker objective (serve a hybrid repaired/stale-shard model without
detection via per-shard signatures), attacker knowledge (architecture,
shard boundaries, one pre-repair-signed shard), attacker control
(artifact cache / update path only), attacker budget (no query/byte
budget -- a single well-timed shard substitution), trusted components
(coordinator, loader, runtime, persisted minimum-epoch state), and the
two named out-of-scope capabilities (fully compromised worker;
signature/loader forgery).}

\subsection{Models, Attack, and Repairs}
\todo{PreActResNet18 / CIFAR-10 / BadNets, poison\_rate=0.1
(data/checkpoint\_manifest.md). Two repairs compared: (a) 20-epoch
global clean fine-tune, (b) ANP (BackdoorBench's own defense script --
Chinh, Priority 1, blocks all of Priority 2 per the 21/09 plan).
State exact commands/configs/seeds per Form C of the workbook -- do
not summarize from memory.}

\subsection{Partitioning and Substitution}
\todo{Describe the \(n\)-shard partition scheme actually used (n=3,
n=6, from src/partition/shard.py) and the substitute() harness, whose
correctness was independently verified via full-rollback
reconstruction (Table~\ref{tab:fullrollback}).}

\subsection{Repair-Share Metric (corrected)}
\label{sec:share}
\todo{Define \(r_s\) / share\_s as actually computed AFTER the fix
described in notes/2026-09-21\_hung\_supervisor.md Section 2: only
learnable parameters, buffers such as \texttt{num\_batches\_tracked},
\texttt{running\_mean}, \texttt{running\_var} excluded. State plainly
that an earlier version of this metric was contaminated by BatchNorm
counters and has been withdrawn -- do not reuse the 19/09 ``25\% per
block'' figures anywhere in this paper.}

\subsection{Metrics}
\todo{ASR and CA definitions and denominators, copied from Project
Agreement.md Section 4 ("Primary outcome and denominator"). State the
Day-5 gate thresholds that were evaluated (ASR $\ge$ 80\% infected,
$\le$ 10\% repaired, $\ge$ 50\% or $\ge$ 30pp for single-shard
recovery) for context, even though the single-shard gate result is
reported as a negative finding, not hidden.}

% =====================================================================
\section{Results}
\label{sec:results}
% Owner: Dan, 24/09, EXCEPT Section~\ref{sec:serving} (Deployment and
% Serving Comparison), owned by Dat. Two rules apply (21/09 plan,
% Section 6): report every k / every strategy / every seed that was
% run, and do not claim a prediction that was not actually made in
% advance.

\subsection{Baseline Reproduction}
\todo{Table~\ref{tab:baseline}: infected ASR/CA, repaired ASR/CA.
Numbers already exist: ASR 97.43\%, CA 89.41\% (infected); ASR 0.90\%,
CA 93.35\% (repaired, 20-epoch fine-tune) -- CONFIRM against the
corrected eval\_baseline.py output before freezing (second-reader note
flagged the earlier hard-coded CA constant; re-verify it was fixed
correctly).}
\begin{table}[h]
\centering
\caption{Baseline reproduction (placeholder -- confirm before freeze)}
\label{tab:baseline}
\begin{tabular}{lcc}
\toprule
Model & ASR (\%) & CA (\%) \\
\midrule
Infected & \todo{97.43} & \todo{89.41} \\
Repaired (20-epoch FT) & \todo{0.90} & \todo{93.35} \\
\bottomrule
\end{tabular}
\end{table}

\subsection{Sanity Checks: Hybrid Models Are Not Broken}
\todo{Table~\ref{tab:fullrollback}: full-rollback reconstruction
(exact match to infected ASR/CA at n=3 and n=6) and hybrid-model CA
range (84.58--93.26\%, from notes/2026-09-19\_dan.md Sections 2.1--2.2)
-- this is what justifies treating the low single-shard ASR as a real
negative result rather than a broken harness.}
\begin{table}[h]
\centering
\caption{Full-rollback verification and hybrid CA range (placeholder)}
\label{tab:fullrollback}
\begin{tabular}{lccc}
\toprule
Config & Rollback ASR (\%) & Rollback CA (\%) & Hybrid CA range (\%) \\
\midrule
n=3 & \todo{97.43} & \todo{89.41} & \todo{84.58--90.98} \\
n=6 & \todo{97.43} & \todo{89.41} & \todo{86.03--93.26} \\
\bottomrule
\end{tabular}
\end{table}

\subsection{Single-Shard Rollback (Negative Result)}
\todo{State plainly: best single-shard rollback ASR = 17.46\%, below
the 50\% gate. Report per-shard breakdown, not just the max. This is
the Day-5 finding -- present it as the attack's lower bound (Section
4 framing of the 21/09 plan), not as a failure to be minimized or an
inflated threat.}

\subsection{\(k\)-Shard Recovery Curve}
\label{sec:kcurve}
\todo{Figure~\ref{fig:kcurve}: ASR and CA vs.\ \(k\) for \(k=1,\dots,6\)
at n=6, from Dan's Priority-1 run. Known endpoints: \(k=0\): ASR
0.90\%; \(k=6\): ASR 97.43\%. Report the full curve, not just where it
crosses 50\% -- per the ``report everything'' rule.}
\begin{figure}[h]
\centering
\todo{\rule{0.9\linewidth}{4cm}}
\caption{ASR and clean accuracy vs.\ number of stale shards $k$
(n=6). Placeholder -- replace with the actual plot from
results/raw/ once Dan's Priority-1 run completes.}
\label{fig:kcurve}
\end{figure}

\subsection{Shard-Selection Strategy Comparison}
\label{sec:strategy}
\todo{Figure/Table~\ref{fig:strategy}: three curves/columns at every
\(k\) -- repair-aware (by corrected share\_s, Section~\ref{sec:share}),
architectural (index order), random (mean of 5 draws). State clearly
whether repair-aware selection reaches the recovery threshold at a
smaller \(k\) than the other two -- report the comparison as run, not
as a confirmed hypothesis (Section 6 of the 21/09 plan).}
\begin{figure}[h]
\centering
\todo{\rule{0.9\linewidth}{4cm}}
\caption{ASR recovery vs.\ $k$ under three shard-selection strategies:
repair-aware, architectural, random (mean $\pm$ range over 5 draws).
Placeholder.}
\label{fig:strategy}
\end{figure}

\subsection{Contrast with a Local Repair (ANP)}
\todo{Repeat Sections~\ref{sec:kcurve}--\ref{sec:strategy} against the
ANP-repaired checkpoint (Chinh's Priority-1 deliverable). Expected
(not yet confirmed) direction: ANP is a local repair and should be
more fragile to rollback than the global 20-epoch fine-tune -- state
this as the motivating question, and report the actual result
regardless of which way it goes.}

\subsection{Repair Locality vs.\ Recovered ASR}
\todo{Figure: corrected locality metric (Section~\ref{sec:share})
plotted against recovered ASR, for both repairs, all shards (Chinh,
Priority 2). This is the figure meant to explain \emph{why} the
attack works when it does.}

\subsection{Second-Seed Replication}
\todo{Table: best configuration re-run on a second seed (Dan,
Priority 2), to show the headline number is not a single-seed
artifact.}

\subsection{Deployment and Serving Comparison}
\label{sec:serving}
% Owner: Ho Du Tuan Dat -- this is YOUR subsection (Hung's own wording:
% "the deployment and serving subsection"). Draft prose here on 24/09;
% the underlying experiment is Priority 2, due 23/09.
\todo{One paragraph + one table, per the 21/09 plan (``Do not build a
system''). Describe: two processes loading a model from a manifest;
comparison of per-shard signature checking vs.\ epoch-manifest
pinning. Table must report exactly these four numbers:
\begin{itemize}
  \item manifest size (bytes)
  \item verification time per load
  \item rollout recovery time after an interrupted update
  \item whether a skewed (mixed-epoch) model is accepted or rejected
\end{itemize}}
\begin{table}[h]
\centering
\caption{Per-shard signature checking vs.\ epoch-manifest pinning (placeholder)}
\label{tab:serving}
\begin{tabular}{lcc}
\toprule
Metric & Per-shard sig.\ check & Epoch-manifest pinning \\
\midrule
Manifest size (bytes) & \todo{--} & \todo{--} \\
Verify time / load & \todo{--} & \todo{--} \\
Rollout recovery time & \todo{--} & \todo{--} \\
Skewed model accepted? & \todo{Yes (undetected)} & \todo{No (rejected)} \\
\bottomrule
\end{tabular}
\end{table}

% =====================================================================
\section{Discussion}
\label{sec:discussion}
\todo{Tie Section~\ref{sec:results} together: single-shard rollback is
weak but silent; multi-shard, repair-aware rollback is the real attack
surface; epoch-manifest pinning closes the gap without a new protocol
(TUF/Uptane already define it -- this paper measures the ML-specific
consequence of its absence, per Project Agreement.md Section 2's
closest-prior-work comparison).}

\subsection{Limitations}
% Folded into Discussion (not its own \section) to protect the hard
% 6-page budget Hung stressed repeatedly. Keep this to one short
% paragraph, not a bulleted checklist -- there is no page budget for it.
\todo{One paragraph, per Form J of the workbook: single architecture
(PreActResNet18), single dataset (CIFAR-10), attack (BadNets) and at
most one more trigger if run, two repair methods only, out-of-scope
threat capabilities as listed in Section~\ref{sec:threat} (fully
compromised worker, signature forgery). Name anything from the
``Not doing'' list in the 21/09 plan (fine-pruning, NAD, 1-epoch
fine-tune, other partitions/architectures/datasets) as explicit future
work in one clause, not a separate list.}

% =====================================================================
\section{Related Work}
\label{sec:related}
% Owner: Chinh, 25/09. Source: Project Agreement.md, "Closest prior
% work and exact remaining difference" -- six items, reuse the stated
% differences verbatim rather than re-deriving them under deadline
% pressure.
\todo{Cover, with the exact differences already written in Project
Agreement.md: Min et al.\ (NeurIPS 2024) on backdoor reactivation via
retuning/querying; DeferBad on engineered unlearning fragility to
forward benign updates; Subnet Replacement Attack on arbitrary
parameter substitution; BadMerging/MergeBackdoor on cross-source model
merging; AIRS on per-shard hash/signature and loader integrity;
TUF/Uptane on rollback and mix-and-match detection at the
snapshot-metadata level.}

% =====================================================================
\section{Conclusion}
\label{sec:conclusion}
% Owner: unassigned in the 21/09 plan ("Someone writes the
% Conclusion") -- confirm who on 25/09.
\todo{Restate the claim from Section 4 of the 21/09 plan in 3--4
sentences: what was measured, the headline numbers (once frozen), and
the practical takeaway (epoch-manifest pinning as a cheap, existing
control). No new numbers here.}

% =====================================================================
\section*{Acknowledgment}
\todo{Standard RIVF acknowledgment line, if applicable -- confirm
requirement with Hung. Cut this first if the paper runs over 6 pages;
it is the lowest-cost section to drop.}

% =====================================================================
\bibliographystyle{IEEEtran}
\bibliography{references}

\end{document}
